% C.3ae: full proof independently audited by /root and /root/frobenius.
\paragraph{C.3ae (Refuted, with an explicit counterexample): ordinary $J^2$ restriction is not surjective in every dimension at least seven.}
For every $N\ge7$ we specify an excellent normal strictly
henselian local complex ring $R$ of dimension $N$, with
regular puncture $X$ and regular punctured principal divisor
$Y$, for which there is a surjection
\begin{equation}\label{eq:C3ae-cokernel}
 \operatorname{coker}\!\left(
 \mathrm{CH}^2(X)_{\mathrm{alg}}
       \longrightarrow\mathrm{CH}^2(Y)_{\mathrm{alg}}\right)
       \longrightarrow \operatorname{Pic}^0(C)(\mathbb C),
 \qquad C:\ z_0^4+z_1^4+z_2^4=0.
\end{equation}
Here the algebraically trivial groups use actual families on
the fixed scheme times a smooth connected parameter curve.
Consequently the ordinary rational map $J^2(X)\to J^2(Y)$
is not surjective. The argument concerns these ordinary
groups; it does not assert that every cycle on $X$ is
annihilated by the invariant constructed below.

\emph{1. The coefficients and the algebraic local scheme.}
Put $m=N-4\ge3$, and let
\[
 W=V(w_0^4+w_1^4+w_2^4+w_3^4)\subset\mathbb P^3,
 \qquad S=W\times C,
 \qquad M=\mathcal O_W(1)\boxtimes K_C.
\]
Write $A_i=(a_{ijk})_{0\le j\le3,\,0\le k\le2}$ for
the following four integer matrices:
\[
\begin{aligned}
 A_0&=\begin{pmatrix}-3&-5&5\\-3&-5&-2\\-5&3&-3\\5&1&-3\end{pmatrix},&
 A_1&=\begin{pmatrix}-4&4&-1\\3&5&2\\-5&0&-1\\0&2&-5\end{pmatrix},\\[2mm]
 A_2&=\begin{pmatrix}0&-1&-1\\0&-2&4\\-1&1&-4\\-5&-1&-3\end{pmatrix},&
 A_3&=\begin{pmatrix}0&-4&1\\-5&-4&-4\\-1&5&1\\4&2&5\end{pmatrix}.
\end{aligned}
\]
Define
\[
 s_i(w,z)=\sum_{j=0}^3\sum_{k=0}^2a_{ijk}w_jz_k,
 \qquad
 \ell(x,w,z)=\sum_{i=0}^3x_i s_i(w,z),
 \qquad
 G(z)=z_0^2z_1^2+z_0^2z_2^2+z_1^2z_2^2.
\]
For $i>3$ we set $s_i=0$. The four $s_i$ have no
common zero on $S$. This is an exact finite verification:
over $\mathbb F_{101}$, in each of the twelve affine charts
$w_j=z_k=1$, the ideal generated by the two Fermat equations
and the four $s_i$ has Groebner basis $[1]$. The common-zero
scheme is projective over $\mathbb Z$, so its image in
$\operatorname{Spec}\mathbb Z$ is closed. A nonempty generic
fibre would force that image to contain every prime, contrary
to the computation at $101$. Thus this particular integer
tuple, without a generic-coefficient qualification, is
base point free over $\mathbb C$.

Set
\[
 B=\mathbb P^m_x\times W\times\mathbb P^2_z,
 \qquad
 \mathcal N=\mathcal O_{\mathbb P^m}(-1)
       \boxtimes\mathcal O_W(-1)\boxtimes\mathcal O_{\mathbb P^2}(-1),
 \qquad T_B=\operatorname{Tot}(\mathcal N).
\]
If $t$ is the tautological fibre vector, define the hypersurface
\[
                   Z=V(F+t\ell G)\subset T_B,
                   \qquad F=z_0^4+z_1^4+z_2^4.
\]
Its zero-section divisor is
\[
 V=\mathbb P^m\times W\times C,
 \qquad N_{V/Z}=L^{-1},
 \qquad L=\mathcal O_{\mathbb P^m}(1)\boxtimes M.
\]
The expression defining $Z$ is a section of the pullback of
$\mathcal O_B(0,0,4)$, so the definition is independent of
trivializations. The nonzero differential $dF$ makes $Z$
smooth near $V$. On $\ell G\ne0$ it is the graph
$t=-F/(\ell G)$. The locus $F=\ell G=0$ has codimension two
in $B$ and produces no additional component of the
hypersurface. It is a Cartier hypersurface in a regular
scheme, hence Cohen--Macaulay, and it is generically
reduced on the graph. It follows that $Z$ is integral.

Contract the zero section of $T_B$ by the Segre coordinates
\[
                         y_{ijk}=t x_iw_jz_k.
\]
This is a proper morphism to the affine cone of the indicated
Segre embedding, and it is an isomorphism off the zero
section. Let $Q$ be the reduced closed image of $Z$, let
$Q^{\mathrm n}$ be its finite normalization, and let
$v^{\mathrm n}$ be the point over the cone vertex. This
point is unique: after shrinking around the vertex, the
smooth proper birational model $Z$ factors through the
normalization and is surjective; its connected exceptional
fibre $V$ cannot map onto more than one point of the finite
normalization fibre. Define
\[
 R=\mathcal O_{Q^{\mathrm n},v^{\mathrm n}}^{\mathrm{sh}}.
\]
This specifies $R$ by the displayed integer equations,
closed image, normalization, and strict henselization.
As in C.3m, properness allows us to shrink only to prove
regularity of the model; this shrinking makes no change
to $R$. Write $\mathcal Z\to\operatorname{Spec}R$ for
the resulting projective regular model. Its exceptional fibre
is $V$, and $X=\mathcal Z\setminus V$ is the regular
puncture of $\operatorname{Spec}R$. The ring has dimension
$m+4=N$ and is excellent and normal. Pointed etale base
change and its limit preserve these assertions and the
first infinitesimal neighbourhood of $V$. More explicitly,
the $y_{ijk}$ generate $(t)$ on each Segre chart, while
every extra normalization function in the maximal ideal
vanishes on the reduced fibre. Thus
$\mathfrak m_R\mathcal O_{\mathcal Z}=\mathcal I_V$,
as in C.3m; the quotients by the squares of these ideals
are unchanged by a pointed etale extension with the same
residue field.

The nonzero nonunit
\[
                      f=\sum_{i,j,k}a_{ijk}y_{ijk}
\]
pulls back to $t\ell$. Its strict transform is
$D=V(F,\ell)\subset T_B$, the total space of $L_H^{-1}$
over
\[
 H=V(\ell)\subset V,
 \qquad L_H=L|_H.
\]
The surjection $\mathcal O_S^{m+1}\to M$ with entries
$s_i$ has a rank-$m$ vector-bundle kernel $E$, and
\[
             H=\mathbb P_{\mathrm{lines}}(E),
             \qquad c_1(L_H)=h+a,
             \qquad a=c_1(M),\quad h=c_1(\mathcal O_H(1)).
\]
Thus $H$ and $D$ are smooth. Their corresponding base
change $\mathcal D\subset\mathcal Z$ is regular, and
$Y=\mathcal D\setminus H=X\cap V(f)$ is the required
regular punctured divisor. The divisors $V$ and
$\mathcal D$ intersect transversely in $H$.

\emph{2. A cycle invariant, with its family compatibility.}
Fix $\eta\in\overline{\mathbb Q}$ with $\eta^4=-1$ and
the line
\[
 D_W=V(w_0-\eta w_1,\,w_2-\eta w_3)\subset W.
\]
Writing $H_W=c_1(\mathcal O_W(1))$, put
$\tau=4[D_W]-H_W$. Adjunction and the degree of the line
give
\begin{equation}\label{eq:C3ae-intersections}
 H_W^2=4,\qquad H_W[D_W]=1,\qquad [D_W]^2=-2,
 \qquad \tau H_W=0,\qquad \tau[D_W]=-9.
\end{equation}

For a smooth projective parameter $P$, Chow localization
and refined pullback along $H\times P\subset\mathcal D\times P$
define a homomorphism
\[
 \sigma_P:\mathrm{CH}^2(Y\times P)
 \longrightarrow
 \frac{\mathrm{CH}^2(H\times P)}
 {(h+a)\mathrm{CH}^1(H\times P)}.
\]
Indeed we lift a class to $\mathrm{CH}^2(\mathcal D\times P)$
and pull back. Two lifts differ by the pushforward of a
class in $\mathrm{CH}^1(H\times P)$, and the
self-intersection formula makes their pullbacks differ
by $c_1(N_{H/\mathcal D})=-(h+a)$ times that class.
This proves independence, including for rational-equivalence
witnesses. These are the same localization and
self-intersection operations used in C.3e.

Since $\operatorname{rank}E=m\ge3$, the first
projective-bundle relation has degree at least three.
Consequently evaluation at $h=-a$ gives an integral
isomorphism
\begin{equation}\label{eq:C3ae-evaluation}
 \frac{\mathrm{CH}^2(H\times P)}
 {(h+a)\mathrm{CH}^1(H\times P)}
                   \simeq \mathrm{CH}^2(S\times P).
\end{equation}
Intersect the right side with $\operatorname{pr}_W^*\tau$
and push forward along $W$. We obtain
\begin{equation}\label{eq:C3ae-rho}
 \rho_P:\mathrm{CH}^2(Y\times P)
       \longrightarrow \mathrm{CH}^1(C\times P)
       =\operatorname{Pic}(C\times P).
\end{equation}
All the operations commute with refined pullback to a
point of $P$. For example, a lift chosen before taking
the fibre is still a valid lift after taking it, and
refined pullbacks commute. Hence $\rho=\rho_{\operatorname{Spec}\mathbb C}$
sends ordinary algebraic-equivalence differences to
$\operatorname{Pic}^0(C)(\mathbb C)$.

If $\Gamma\in\mathrm{CH}^2(X\times P)$ and
$\widetilde\Gamma$ is a closure on $\mathcal Z\times P$,
write its refined restriction to $V\times P$ as
\[
 b=b_0+h b_1+h^2 b_2,
 \qquad b_i\in\mathrm{CH}^{2-i}(S\times P).
\]
Transverse refined pullback to $H$ and
\eqref{eq:C3ae-evaluation} give the exact formula
\begin{equation}\label{eq:C3ae-source-invariant}
 \rho_P(\Gamma|_{Y\times P})
   =\operatorname{pr}_{C\times P,*}
            \bigl(\tau(b_0-a b_1+a^2b_2)\bigr).
\end{equation}
There is no presumed extension of a cycle over a varying
etale neighbourhood in this formula: the parameter $P$
is fixed, and the closure is on the fixed proper model
$\mathcal Z\times P$.

\emph{3. The first-neighbourhood obstruction for a cycle class.}
We give the required Chern-class argument explicitly.
All differential sheaves in this argument are relative
to $\mathbb C$.
Let $V_2\subset\mathcal Z$ be defined by $\mathcal I_V^2$.
Its square-zero ideal is $L$, and its class is
$\kappa\in H^1(V,T_V\otimes L)$. For any smooth
projective $P$, the conormal sequence of $V\times P$
in $\mathcal Z\times P$ and its second exterior power are
\[
\begin{gathered}
 0\longrightarrow L\longrightarrow
 \Omega^1_{\mathcal Z\times P}|_{V\times P}
 \longrightarrow\Omega^1_{V\times P}\longrightarrow0,\\
 0\longrightarrow L\otimes\Omega^1_{V\times P}
 \longrightarrow\Omega^2_{\mathcal Z\times P}|_{V\times P}
 \longrightarrow\Omega^2_{V\times P}\longrightarrow0.
\end{gathered}
\]
The second sequence is exact because the kernel of
the first has rank one. Its cohomology boundary is
contraction with $\kappa$.
This can be checked by taking local splittings of the
first sequence: their differences are derivations
representing $\kappa$, and the induced difference on
a two-form is contraction in either factor with the
usual alternating sign.

For a vector bundle on $\mathcal Z\times P$, its Hodge
Chern character restricts through the middle term of
the second sequence. Concretely, on a trivializing cover
the matrices $g_{uv}^{-1}dg_{uv}$ represent the Atiyah
class, and $\tfrac12\operatorname{tr}(\operatorname{At}^2)$
represents its degree-two Chern character. These matrices
are defined before restriction, so their trace cup
products give the asserted lift. The formula is
compatible with direct sums and exact sequences; on a
flag bundle it reduces to $\tfrac12\sum c_1(L_i)^2$,
which identifies it with the usual Hodge realization
of the Chow Chern character. Thus it also applies to
virtual vector bundles and rational combinations.

For our closure class we may choose such a rational
virtual bundle whose degree-two Chow Chern character
is $\widetilde\Gamma$. Precisely, the regular scheme
$\mathcal Z\times P$ is finite dimensional, has affine
diagonal and an ample invertible sheaf, and the rational
Chern character is an isomorphism onto its Chow group
\citep[Lemma 42.58.1, Tag 0FEY]{cyStacksChernObstruction}.
One may equally descend this finite collection of cycles,
bundles and maps to a pointed etale neighbourhood and
work on a smooth finite-type complex model; its restriction
to $V_2\times P$ is the same. Naturality of Chern classes
under refined pullback
\citep[Remark 42.59.9, Tag 0FF8]{cyStacksChernObstruction}
therefore gives
\begin{equation}\label{eq:C3ae-obstruction-zero}
 \kappa\mathbin{\lrcorner}\operatorname{cl}_{\mathrm{Hdg}}(b)
       =0
 \quad\hbox{in }H^3(V\times P,L\otimes\Omega^1_{V\times P}).
\end{equation}
Only this necessary vanishing is used. We make no
converse assertion about deforming cycles from their
Hodge classes.

The class $\kappa$ has only a curve-tangent component,
since the construction has a morphism to
$\mathbb P^m\times W$. The normal sequence calculation
from C.3m gives
\begin{equation}\label{eq:C3ae-kappa}
 \kappa=\sum_{i,j,k}a_{ijk}x_iw_j\otimes\delta(z_kG),
 \qquad
 \delta:H^0(C,\mathcal O_C(5))\longrightarrow H^1(C,\mathcal O_C).
\end{equation}
Up to the same harmless overall sign, the three values
$\delta(z_kG)$ are a basis. In the Fermat Jacobian ring
they are respectively
$z_0z_1^2z_2^2$, $z_0^2z_1z_2^2$, and
$z_0^2z_1^2z_2$. The source basis of $H^0(C,K_C)$
used below is the Serre-dual basis to these three
values; it need not be identified with a chosen
coordinate basis of that space.

\emph{4. The exact rank calculation and the Hodge projection.}
Put
\[
 \mathcal R=\mathbb C[w_0,w_1,w_2,w_3]/(w_0^3,w_1^3,w_2^3),
 \qquad \dim\mathcal R_3=16,\quad\dim\mathcal R_4=19.
\]
The Euler and normal sequences on the quartic give
\[
 H^1(W,T_W(-1))\simeq\mathcal R_3,
 \qquad
 H^1(W,T_W)_{\mathrm{pol}}\simeq\mathcal R_4.
\]
For completeness, $H^1(T_{\mathbb P^3}(-1)|_W)=0$,
and quotienting $H^0(\mathcal O_W(3))$ by the four
partial derivatives gives the first equality. In degree
zero, quotienting $H^0(\mathcal O_W(4))$ by linear
vector fields gives the $19$-dimensional kernel of
\[
                   H^1(T_W)\longrightarrow H^1(T_{\mathbb P^3}|_W).
\]
The latter
target has dimension one by the Euler sequence.
The K3 equalities $K_W\simeq\mathcal O_W$ and
$h^{1,1}(W)=20$ identify this kernel with the
polarized deformation space. Its Serre-dual annihilator
in $H^1(W,\Omega^1_W)$ is exactly $\mathbb C H_W$:
the polarization extends to every embedded deformation,
and its class is nonzero since $H_W^2=4$.

Serre duality makes
\[
 \frac{H^1(W,\Omega^1_W)}{\mathbb C H_W}
      \simeq\mathcal R_4^*,
 \qquad H^1(W,\Omega^1_W(1))\simeq\mathcal R_3^*.
\]
Multiplication by $w_j$ on these Hodge groups is dual
to multiplication $w_j:\mathcal R_3\to\mathcal R_4$.
In particular every such multiplication kills $H_W$.
The displayed identifications use dual vector spaces,
not an unmentioned self-dual monomial basis.

Order the degree-three and degree-four monomials
lexicographically by their exponent tuples in
$\{0,1,2\}^4$. In the dual bases, the map
\begin{equation}\label{eq:C3ae-rank-map}
 \mathcal R_4^*\otimes\mathbb C^3
       \longrightarrow\mathbb C^4\otimes\mathcal R_3^*
\end{equation}
induced by \eqref{eq:C3ae-kappa} has rows $(i,u)$,
columns $(k,v)$, and entry
\[
                  M_{(i,u),(k,v)}
                     =\sum_{\{j:\,u+e_j=v\}} a_{ijk}.
\]
It is a $64$ by $57$ integer matrix. The determinant of
the submatrix consisting of its first $57$ rows is
$47$ modulo $101$. Hence \eqref{eq:C3ae-rank-map} is
injective over $\mathbb C$. The retained code records
the monomial order, matrices and exact finite-field
elimination, and independent implementations reproduce
both this minor and full column rank.

Now project \eqref{eq:C3ae-obstruction-zero} to
\begin{equation}\label{eq:C3ae-target-component}
 \begin{split}
 H^0(\mathbb P^m,\mathcal O(1))
  \otimes H^1(W,\Omega^1_W(1))
  \otimes H^1(C,K_C)
  \otimes H^1(P,\mathcal O_P).
 \end{split}
\end{equation}
Because $\kappa$ contracts only a curve differential,
the relevant source component of $\operatorname{cl}(b_0)$ is
\begin{equation}\label{eq:C3ae-source-component}
 H^1(W,\Omega^1_W)\otimes H^0(C,K_C)
                          \otimes H^1(P,\mathcal O_P).
\end{equation}
Terms with a projective-space differential cannot
contribute to \eqref{eq:C3ae-target-component}. They
also vanish by the Euler-sequence equality
\[
                       H^1(\mathbb P^m,\Omega^1(1))=0.
\]
The action on
\eqref{eq:C3ae-source-component}, modulo $\mathbb C H_W$
in the first factor, is exactly the injective map
\eqref{eq:C3ae-rank-map} tensored with
$H^1(P,\mathcal O_P)$. Thus the primitive part of
\eqref{eq:C3ae-source-component} is zero.

Pairing its first factor with $\tau$ annihilates the
remaining $H_W$ part by \eqref{eq:C3ae-intersections}.
It follows that the line bundle
$\operatorname{pr}_{C\times P,*}(\tau b_0)$ has zero
mixed component in $H^0(C,K_C)\otimes H^1(P,\mathcal O_P)$.
The correction $-a b_1+a^2b_2$ in
\eqref{eq:C3ae-source-invariant} has no effect on this
conclusion. Indeed $H^1(W,\mathcal O_W)=0$ and
$H^0(W,\Omega_W^1)=0$, so the only varying mixed
component of a divisor on $S\times P$ comes from
$C\times P$. Multiplication by $a$ supplies the
$H_W$ factor, which pairs to zero with $\tau$.
All other divisor components are constant in $P$;
$b_2$ has codimension zero and is constant as well.
This verifies the boundary correction rather than
assuming that a chosen lift is unique.

\emph{5. Actual ordinary families and the nonzero target.}
Every generator of $\mathrm{CH}^2(X)_{\mathrm{alg}}$ is
a difference of fibres of a cycle on $X$ times a
smooth connected curve. Compactify the curve smoothly
and close that cycle in the product. The fibres at
the original points are unchanged. Thus it suffices
to take the parameter $P$ above to be a smooth
projective connected curve and a fixed actual cycle
on $X\times P$.

The line bundle in \eqref{eq:C3ae-source-invariant}
defines, after subtracting its value at a base point,
a morphism $P\to\operatorname{Pic}^0(C)$. Its mixed
Hodge component is zero by step 4. By the universal
property of $J(P)$, this morphism factors through a
homomorphism $J(P)\to\operatorname{Pic}^0(C)$, whose
map on $H^1(\mathcal O)$ is exactly that mixed
component, as in the Poincare calculation in C.3m
\citep[Theorem 1.2 and Lemma 9.5]{D_MilneJacobians2021}.
A homomorphism of complex abelian varieties
with zero such map is zero: complex conjugation gives
zero on holomorphic differentials, hence zero
differential, and a homomorphism with zero differential
is constant. The classifying morphism is therefore
constant. Fibre compatibility of $\rho_P$ proves
\begin{equation}\label{eq:C3ae-source-zero}
                \rho\bigl(\mathrm{CH}^2(X)_{\mathrm{alg}}|_Y\bigr)=0.
\end{equation}
This proof uses fixed proper parameter curves. It
does not require extending a class from a varying
etale neighbourhood across a compactified parameter.

Pulling the divisors $D_W$ and a divisor on $C$ to
$Y$ gives a homomorphism
\[
 \theta:\operatorname{Pic}^0(C)(\mathbb C)
               \longrightarrow \mathrm{CH}^2(Y)_{\mathrm{alg}}.
\]
Its image is ordinarily algebraically trivial,
as is seen from the Poincare line bundle on
$C\times\operatorname{Pic}^0(C)$: two points of that
connected projective abelian variety can be joined
by a chain of smooth connected curves, and we restrict
the actual family to those curves. Taking the corresponding
extension on $\mathcal D$ and applying the definition
of $\rho$ yields the integral identity
\[
                         \rho\theta=-9\,\mathrm{id}.
\]
Multiplication by $9$ is surjective on the complex
points of an abelian variety. Equations
\eqref{eq:C3ae-source-zero} and the last identity prove
\eqref{eq:C3ae-cokernel}.

For one specified nontorsion witness, choose
$u\in\mathbb C$ transcendental over
$\overline{\mathbb Q}$, choose $v$ with
$v^4=-1-u^4$, and set
\[
                       p=(u:v:1),\qquad o=(1:\eta:0).
\]
Then $[p-o]\in\operatorname{Pic}^0(C)(\mathbb C)$ is
nontorsion: otherwise its Abel point would be a
torsion point defined over $\overline{\mathbb Q}$,
and the Abel closed immersion, defined over that
field, would force $p$ to be algebraic
\citep[Proposition 2.3]{D_MilneJacobians2021}. The target
cycle $\theta([p-o])$ therefore has nonzero image
$-9[p-o]$ under $\rho$ and cannot lie in the ordinary
source image, even after tensoring with $\mathbb Q$.

\emph{Finite certificates and proof scope.}
The rank calculation took $0.048689819$ seconds and
the twelve-chart base point calculation took
$4.204962002$ seconds in the retained implementations.
Independent checks used different matrix and Groebner
basis algorithms; in particular the $57$ by $57$
minor was computed independently. Those computations
certify the stated finite linear algebra and absence
of base points. The localization, Chern obstruction,
boundary correction and proper-family argument above
are the separate reasons that these data give the
claimed ordinary Chow-group counterexample.
